
\documentclass[11pt,a4paper]{article}

% Encodage et langue
\usepackage[T1]{fontenc}
\usepackage[utf8]{inputenc}
\usepackage[french]{babel}
\usepackage{lmodern}

% Mise en page et mathématiques
\usepackage[a4paper,margin=2.5cm]{geometry}
\usepackage{amsmath,amssymb,mathtools}
\usepackage{booktabs}
\usepackage{enumitem}
\usepackage{xcolor}
\usepackage{hyperref}

\hypersetup{
    colorlinks=true,
    linkcolor=blue!50!black,
    urlcolor=blue!50!black
}

\title{\textbf{Détection de contours par la méthode de Sobel}\\
       \large Fondements mathématiques et description de l'algorithme}
\author{Ronhy RAKOTONDRAFARA}
\date{}

\begin{document}
\maketitle

\section{Introduction}

La méthode de Sobel est un opérateur de détection de contours en traitement
d'images. Elle permet de repérer les zones où l'intensité lumineuse varie
rapidement. Ces variations correspondent généralement aux frontières entre
des régions de l'image.

L'idée mathématique consiste à approximer les dérivées spatiales de l'image
selon deux directions perpendiculaires : l'horizontale et la verticale.
À partir de ces deux dérivées, on calcule l'intensité et la direction du
gradient, puis on peut éventuellement appliquer un seuil pour obtenir une
image binaire des contours.

\section{Représentation mathématique de l'image}

Une image en niveaux de gris est représentée par une fonction discrète
\[
    I : \{0,\ldots,H-1\}\times\{0,\ldots,W-1\}
    \longrightarrow \mathbb{R},
\]
où $H$ est la hauteur de l'image, $W$ sa largeur, et $I(y,x)$ l'intensité
du pixel situé à la ligne $y$ et à la colonne $x$.

Dans une image codée sur 8 bits, l'intensité est généralement comprise entre
$0$ et $255$ :
\[
    0 \leq I(y,x) \leq 255.
\]
La valeur $0$ représente le noir et la valeur $255$ le blanc.

\subsection{Conversion en niveaux de gris}

Lorsqu'une image est en couleur, on peut construire une intensité unique à
partir des composantes rouge $R$, verte $G$ et bleue $B$ :
\[
    I(y,x)=0.299R(y,x)+0.587G(y,x)+0.114B(y,x).
\]
Les coefficients donnent davantage de poids à la composante verte, à laquelle
l'œil humain est généralement plus sensible.

Si l'image est déjà en niveaux de gris, cette étape n'est pas nécessaire.

\section{Principe du gradient}

Pour une image continue $I(x,y)$, le gradient est le vecteur formé par les
dérivées partielles de l'intensité :
\[
    \nabla I(x,y)=
    \begin{pmatrix}
        \dfrac{\partial I}{\partial x}\\[4pt]
        \dfrac{\partial I}{\partial y}
    \end{pmatrix}.
\]

Les deux composantes du gradient sont :
\[
    G_x=\frac{\partial I}{\partial x},
    \qquad
    G_y=\frac{\partial I}{\partial y}.
\]

\begin{itemize}[leftmargin=1.2cm]
    \item $G_x$ mesure la variation de l'intensité dans la direction
    horizontale.
    \item $G_y$ mesure la variation de l'intensité dans la direction
    verticale.
\end{itemize}

Une valeur élevée de l'une de ces composantes indique une variation
importante de luminosité dans la direction correspondante. Les contours
sont donc associés aux zones où la norme du gradient est élevée.

\section{Masques de Sobel}

Dans une image numérique, les dérivées sont approximées à l'aide de masques
de convolution de taille $3\times3$.

Le masque horizontal est :
\[
    K_x=
    \begin{bmatrix}
        -1 & 0 & 1\\
        -2 & 0 & 2\\
        -1 & 0 & 1
    \end{bmatrix}.
\]

Le masque vertical est :
\[
    K_y=
    \begin{bmatrix}
        -1 & -2 & -1\\
         0 &  0 &  0\\
         1 &  2 &  1
    \end{bmatrix}.
\]

Ces masques associent deux opérations :
\begin{enumerate}[leftmargin=1.2cm]
    \item une différence entre les intensités de part et d'autre du pixel
    central, afin d'estimer une dérivée ;
    \item une pondération des pixels voisins, qui donne un effet de
    lissage et rend l'estimation moins sensible aux petites variations.
\end{enumerate}

\section{Convolution et calcul des gradients}

Soit $K$ un masque de taille $3\times3$. La convolution discrète de l'image
$I$ par ce masque est définie par :
\[
    (I*K)(y,x)=
    \sum_{i=-1}^{1}\sum_{j=-1}^{1}
    I(y-i,x-j)\,K(i,j).
\]

Les indices négatifs ou dépassant les dimensions de l'image apparaissent
lorsqu'on traite les pixels situés sur les bords. Pour conserver les
dimensions de l'image, on prolonge ici l'image en répétant la valeur du
pixel de bord le plus proche. Cette règle est appelée prolongement par
réplication des bords.

Les deux composantes du gradient sont alors calculées par :
\[
    G_x=I*K_x,
    \qquad
    G_y=I*K_y.
\]

Pixel par pixel, on obtient :
\[
    G_x(y,x)=
    \sum_{i=-1}^{1}\sum_{j=-1}^{1}
    I(y-i,x-j)K_x(i,j),
\]
\[
    G_y(y,x)=
    \sum_{i=-1}^{1}\sum_{j=-1}^{1}
    I(y-i,x-j)K_y(i,j).
\]

Les matrices $G_x$ et $G_y$ ont les mêmes dimensions que l'image initiale.

\section{Magnitude du gradient}

La magnitude représente l'intensité globale de la variation lumineuse,
indépendamment de sa direction. Elle est calculée à partir de la norme
euclidienne du gradient :
\[
    M(y,x)=\sqrt{G_x(y,x)^2+G_y(y,x)^2}.
\]

Une magnitude proche de zéro indique une région relativement uniforme.
Une magnitude élevée indique une variation importante de l'intensité et
peut donc correspondre à un contour.

La direction du gradient peut également être exprimée par :
\[
    \theta(y,x)=
    \operatorname{atan2}\bigl(G_y(y,x),G_x(y,x)\bigr).
\]
Cette direction est perpendiculaire à l'orientation locale du contour.

\section{Normalisation de la magnitude}

Les valeurs de $M$ ne sont pas nécessairement comprises entre $0$ et $255$.
Pour obtenir une représentation sur 8 bits, on effectue une normalisation
min-max.

On note :
\[
    M_{\min}=\min_{y,x}M(y,x),
    \qquad
    M_{\max}=\max_{y,x}M(y,x).
\]

Si $M_{\max}\neq M_{\min}$, la magnitude normalisée est :
\[
    M_N(y,x)=
    255\,\frac{M(y,x)-M_{\min}}{M_{\max}-M_{\min}}.
\]

Ainsi :
\[
    0\leq M_N(y,x)\leq255.
\]

Dans le cas particulier où $M_{\max}=M_{\min}$, toutes les valeurs sont
identiques ; l'image normalisée est alors définie comme une image nulle :
\[
    M_N(y,x)=0.
\]

\section{Seuillage facultatif}

Pour obtenir une image binaire des contours, on peut appliquer un seuil
$T$ à la magnitude normalisée. Le résultat $E$ est défini par :
\[
    E(y,x)=
    \begin{cases}
        255, & \text{si } M_N(y,x)\geq T,\\
        0,   & \text{sinon}.
    \end{cases}
\]

Les pixels dont la magnitude atteint ou dépasse le seuil sont considérés
comme des contours. Les autres pixels sont supprimés.

Un seuil élevé conserve principalement les variations les plus fortes,
mais peut éliminer des contours peu contrastés. Un seuil faible conserve
davantage de détails, au risque de laisser apparaître des variations
non pertinentes.

Si aucun seuil n'est appliqué, le résultat reste l'image en niveaux de
gris $M_N$, dans laquelle l'intensité des pixels représente la force du
gradient.

\section{Déroulement mathématique de l'algorithme}

Pour une image d'entrée $I$, la méthode suit les étapes suivantes :
\begin{enumerate}[leftmargin=1.2cm]
    \item \textbf{Préparation :} représenter l'image en niveaux de gris,
    si elle est initialement en couleur.
    \item \textbf{Gradient horizontal :} calculer
    \[
        G_x=I*K_x.
    \]
    \item \textbf{Gradient vertical :} calculer
    \[
        G_y=I*K_y.
    \]
    \item \textbf{Magnitude :} calculer, pour chaque pixel,
    \[
        M=\sqrt{G_x^2+G_y^2}.
    \]
    \item \textbf{Normalisation :} transformer les valeurs de $M$ en
    intensités comprises entre $0$ et $255$, selon la formule min-max.
    \item \textbf{Seuillage (facultatif) :} si un seuil $T$ est choisi,
    produire l'image binaire $E$ à partir de $M_N$.
\end{enumerate}

\section{Interprétation et limites}

La méthode de Sobel est simple et relativement peu coûteuse en calcul.
Elle fournit une estimation du gradient dans deux directions et permet
ainsi de mettre en évidence les transitions d'intensité.

Cependant, le masque de Sobel ne supprime pas complètement le bruit.
Des variations locales non significatives peuvent donc produire des
magnitudes élevées. Un pré-lissage de l'image peut être envisagé lorsque
le bruit est important.

De plus, la normalisation min-max dépend des valeurs minimale et maximale
présentes dans l'image traitée. Le contraste de l'image de magnitude peut
donc varier d'une image à l'autre.

Enfin, Sobel détecte les variations d'intensité : un contour n'est pas
nécessairement une frontière sémantique entre deux objets. Il s'agit avant
tout d'une discontinuité locale de luminosité.

\section{Conclusion}

L'opérateur de Sobel estime les dérivées spatiales d'une image à l'aide de
deux masques $3\times3$. Les composantes $G_x$ et $G_y$ permettent de
calculer la magnitude du gradient, qui mesure la force des variations
d'intensité. Après normalisation, un seuillage facultatif permet de
transformer cette information en une carte binaire des contours.

La chaîne mathématique essentielle de la méthode peut être résumée par :
\[
    I
    \longrightarrow
    (G_x,G_y)
    \longrightarrow
    M=\sqrt{G_x^2+G_y^2}
    \longrightarrow
    M_N
    \longrightarrow
    E\ \text{(si seuillage)}.
\]

\end{document}
